\section{Methodology}
\label{sec:method}

We study reward augmentation for vision-language-action (VLA) policy optimization in LIBERO.
Rather than training a policy from scratch, we initialize from an author-provided supervised fine-tuning (SFT) checkpoint and compare two reward settings under the same initialization, environment, and optimization pipeline:
\begin{enumerate}
    \item a baseline reward using only the environment reward $R_{1,t}$, and
    \item our proposed reward using a combined objective $R_t = R_{1,t} + \lambda R_{2,t}$.
\end{enumerate}
This design isolates the effect of the proposed spatial reward term $R_{2,t}$ while controlling for the backbone model and starting point.

\subsection{Policy Initialization}

Let $\pi_\theta(a_t \mid o_t, q)$ denote a VLA policy that predicts a robot action $a_t$ from the current observation $o_t$ and language instruction $q$.
Instead of learning $\theta$ from scratch, we initialize $\theta$ from an author-provided SFT checkpoint:
\begin{equation}
    \theta_0 \leftarrow \theta_{\text{SFT}}.
\end{equation}
This initialization gives a strong and reproducible starting point, allowing us to focus on whether the additional reward term improves downstream policy optimization.

\subsection{Baseline Environment Reward}

During rollout in LIBERO, the simulator provides an environment reward
\begin{equation}
    R_{1,t} = r_{\text{env}}(s_t, a_t, s_{t+1}),
\end{equation}
where $s_t$ and $s_{t+1}$ denote the environment states before and after action $a_t$.
In LIBERO, this reward is typically sparse and tied to task completion or progress.

If we use only the simulator reward, the optimization objective reduces to the baseline setting
\begin{equation}
    R^{\text{base}}_t = R_{1,t}.
\end{equation}

\subsection{Spatial Reward Augmentation}

Our key idea is to supplement the sparse environment reward with a spatial-consistency reward derived from latent alignment.
Given the current observation $o_t$, the policy produces an internal feature representation $h_t$.
We project it into a policy-side spatial latent:
\begin{equation}
    z^{\text{pol}}_t = g_\phi(h_t),
\end{equation}
where $g_\phi(\cdot)$ is a lightweight learnable projection head.

In parallel, we obtain a reference spatial latent from a frozen spatial encoder $f_{\text{ref}}(\cdot)$:
\begin{equation}
    z^{\text{ref}}_t = f_{\text{ref}}(o_t).
\end{equation}

We define the auxiliary spatial reward $R_{2,t}$ as the cosine similarity between the two latent representations:
\begin{equation}
    R_{2,t} = \mathrm{cos}\!\left(z^{\text{pol}}_t, z^{\text{ref}}_t\right)
    = \frac{\left(z^{\text{pol}}_t\right)^\top z^{\text{ref}}_t}
    {\left\|z^{\text{pol}}_t\right\|_2 \left\|z^{\text{ref}}_t\right\|_2}.
\end{equation}

Intuitively, $R_{2,t}$ encourages the policy to produce latent representations that are more spatially aligned with a frozen spatial reference model, thereby injecting additional structure into policy optimization beyond the sparse task reward alone.

\subsection{Combined Reward Objective}

Our proposed reward combines the original environment signal and the spatial reward:
\begin{equation}
    R_t = R_{1,t} + \lambda R_{2,t},
\end{equation}
where $\lambda \ge 0$ controls the contribution of the spatial reward.

This yields two experimental branches:
\begin{align}
    \text{Baseline:} \quad & R_t = R_{1,t}, \\
    \text{Ours:} \quad & R_t = R_{1,t} + \lambda R_{2,t}.
\end{align}
Because both branches share the same initial checkpoint and rollout environment, the performance gap directly measures the impact of the proposed spatial reward term.

\subsection{Policy Optimization}

Let $\tau = (s_0, a_0, s_1, a_1, \dots)$ denote a rollout trajectory sampled from $\pi_\theta$ in the simulator.
We optimize the policy parameters to maximize the expected return under the chosen reward definition:
\begin{equation}
    J(\theta)
    = \mathbb{E}_{\tau \sim \pi_\theta}
    \left[
        \sum_{t=0}^{T-1} \gamma^t R_t
    \right],
\end{equation}
where $\gamma$ is the discount factor.

Under the baseline setting, $R_t = R_{1,t}$.
Under our method, $R_t = R_{1,t} + \lambda R_{2,t}$.
The corresponding policy-gradient objective can be written as
\begin{equation}
    \nabla_\theta J(\theta)
    =
    \mathbb{E}_{\tau \sim \pi_\theta}
    \left[
        \sum_{t=0}^{T-1}
        \nabla_\theta \log \pi_\theta(a_t \mid o_t, q)\,
        \hat{A}_t
    \right],
\end{equation}
where $\hat{A}_t$ is the advantage estimate computed from the reward sequence.
In our method, $\hat{A}_t$ is induced by the augmented reward and thus reflects both task success and spatial alignment.

\subsection{Experimental Logic}

The methodology is designed to answer a focused question:
\emph{given the same author-provided SFT initialization, does adding the spatial reward term improve policy learning beyond the original environment reward alone?}

Accordingly, our comparison keeps the following fixed:
\begin{itemize}
    \item initial checkpoint,
    \item benchmark environment,
    \item observation and language inputs, and
    \item optimization pipeline.
\end{itemize}
The only intended change is the reward definition.
This makes the comparison between $R_{1,t}$ and $R_{1,t} + \lambda R_{2,t}$ a controlled test of the proposed method.
